\section{Related Work}
\label{sec:related}

AI and machine learning (ML) tools are increasingly being deployed on both sides of healthcare administration. On the insurer side, a recent scoping review identifies eight application domains for AI across health insurance, including claims processing, fraud detection, utilization review, and coverage decision support \citep{ramezani2025ai}. On the provider side, \citet{owolabi2025ml} show that elastic net classifiers (regularized linear models that combine L1 and L2 penalties to handle correlated predictors) can predict which hospital admission denials will be overturned on appeal with 84\% accuracy; after deployment of this system, monthly appeal volumes and the number of overturned denials nearly doubled. However, these systems each optimize for one party in isolation; neither models the strategic interaction dynamics that emerge when both sides deploy AI simultaneously.

Multi-agent simulation using large language models offers a promising methodology for studying such strategic interactions. \citet{fan2025aihospital} introduce AI Hospital, in which simulated doctor, patient, examiner, and chief physician agents interact across multi-turn diagnostic conversations. Their work demonstrates that LLM diagnostic accuracy degrades substantially in interactive settings compared to static question-answering benchmarks (even GPT-4 achieves less than half the upper-bound score) and that a multi-doctor dispute resolution mechanism, where multiple LLM physicians deliberate on a diagnosis, can partially recover this lost accuracy. This finding establishes that multi-agent simulation can surface emergent behaviors not visible in single-agent evaluation, motivating its application to adversarial settings like prior authorization.

Our work connects these two threads, single-sided AI optimization and multi-agent LLM simulation, and extends them into an adversarial setting. The single-agent systems reviewed above each optimize for one side of the PA transaction. In contrast, we model both sides simultaneously within a shared clinical encounter, making the strategic interaction itself the object of study. Unlike collaborative multi-agent settings in which agents share a diagnostic goal \citep{fan2025aihospital}, our agents hold opposing objectives: the provider seeks to maximize approved services, while the insurer seeks to minimize unnecessary coverage. They operate under information asymmetry: the provider receives only the clinical treatment guideline, while the insurer receives only the payer's coverage policy. We further introduce behavioral strategy as an experimental variable by assigning cooperative or adversarial instructions through each agent's system prompt, enabling controlled comparison of strategic postures within the same clinical scenario.
